\documentclass[11pt]{article}
\usepackage[a4paper,margin=1in]{geometry}
\usepackage{amsmath,amssymb,amsthm}
\usepackage{booktabs}
\usepackage{graphicx}
\usepackage{xcolor}
\usepackage{hyperref}
\hypersetup{colorlinks=true,linkcolor=blue,urlcolor=blue,citecolor=blue}
\newcommand{\rt}[1]{\texttt{#1}}
\newcommand{\fn}[1]{\texttt{#1}}
\newcommand{\cmark}{\ensuremath{\checkmark}}
\newcommand{\xmark}{\ensuremath{\times}}
\providecommand{\ket}[1]{|#1\rangle}
\newcommand{\wcc}{\mathrm{wcc}}
\newcommand{\rec}{\mathrm{rec}}
\newcommand{\ifgraphic}[2]{\IfFileExists{#1}{\includegraphics[width=#2]{#1}}%
  {\fbox{\parbox{#2}{\centering\small figure \texttt{\detokenize{#1}} pending}}}}
\newcommand{\iftable}[1]{\IfFileExists{#1}{\input{#1}}%
  {\fbox{\parbox{0.9\textwidth}{\centering\small table
  \texttt{\detokenize{#1}} pending}}}}

\theoremstyle{plain}
\newtheorem{theorem}{Theorem}
\newtheorem{proposition}{Proposition}
\theoremstyle{definition}
\newtheorem{definition}{Definition}

\title{\textbf{Report R10 --- Permutation circuits: the same 256 rules with
$V=X$}\\
\large A control for the sector/attractor machinery, open boundaries}
\author{QCA Sector Analysis project (Tier 1)}
\date{2026-07-30 \quad\small(\fn{xgate\_version} \texttt{1x.1}, \rt{obc0})}

\begin{document}
\maketitle

\begin{abstract}
The Tier-1 machinery separates two decompositions of the computational basis: the
weakly connected components (sectors, exact for both the monitored and the
unmonitored dynamics) and the terminal strongly connected components (monitored
attractors). In the Hadamard family those two disagree strongly, and R9 is largely
about the disagreement. This report runs the identical analysis on the same $256$
symbol tables with the active gate changed from the Hadamard to $X$. Every local
symbol is then a function on bits, the transition graph has out-degree one, and
three things follow at once: the two decompositions become \emph{identical} in
number, sectors and basins coincide, and the $16$ rules built from \rt{I} and
\rt{V} alone are reversible cellular automata whose sectors are exactly the
permutation cycles. We verify all of this over $256$ rules, confirm that the
partition sum rule and the exclusion curve carry over unchanged, and record the
one place where the two families genuinely differ in physics rather than in
bookkeeping: the permutation circuits fragment into exponentially many,
exponentially large sectors whose \emph{attractors are almost always short}. Only
$24$ of $256$ rules have an exponentially long cycle, and for the $80$ V-free
rules the number is zero. Swept to $N=24$ for all $256$ rules. \rt{obc0} only;
\rt{pbc} is held.
\end{abstract}

\tableofcontents

\section{What changes when $V=X$}
\label{sec:setup}

The symbol table is untouched. Each entry still says what happens at a site for
one neighbour pattern; only the active gate is different:

\begin{center}
\begin{tabular}{cll}
\toprule
symbol & Hadamard family (R1--R9) & this report \\
\midrule
\rt{I} & identity & identity \\
\rt{V} & Hadamard & $X$, i.e.\ $x_i\mapsto1-x_i$ \\
\rt{D} & reset to $\ket0$ & reset to $\ket0$ \\
\rt{E} & reset to $\ket1$ & reset to $\ket1$ \\
\bottomrule
\end{tabular}
\end{center}

\noindent All four are now \emph{functions} on the computational basis, so the
circuit is one too, and $\operatorname{succ}(x)$ is a single state. Three
consequences, and they are the reason this family is a useful control.

\begin{proposition}[functional graph]
The one-cycle transition graph has out-degree exactly $1$.
\end{proposition}

\begin{theorem}[F1: one cycle per component]
In a functional graph every weakly connected component contains exactly one
cycle. Hence for every X-gate rule, every $N$ and every boundary condition,
\begin{equation}
n_\rec \;=\; n_\wcc ,
\label{eq:F1}
\end{equation}
the basin of each cycle is its entire component, and no state lies in two basins.
\end{theorem}

\noindent That identity is exactly what fails in the Hadamard family, and fails
badly: R9's rule $22$ at $N=12$ has three monitored attractors, each a single
state, sitting inside two sectors that cover all $4096$ states. There A4 could
only be asserted in one direction. Here the two tiers coincide by construction,
so any disagreement in the code would show up immediately --- which makes this
family a test of the machinery as much as a physics question.

\subsection{What ``reversible'' means here, and what it does not}
\label{sec:reversible-def}

The word is overloaded in this area, so it is fixed explicitly.

\begin{definition}
A rule is \emph{reversible at $(N,\rt{bc})$} if its one-cycle map
$\{0,1\}^N\to\{0,1\}^N$ is a bijection.
\end{definition}

\begin{proposition}
A rule is reversible, for every $N\ge3$ and both boundary conventions, iff its
symbol table uses only \rt{I} and \rt{V}. There are $2^4=16$ such rules:
$51$, $54$, $57$, $60$, $99$, $102$, $105$, $108$, $147$, $150$, $153$, $156$,
$195$, $198$, $201$, $204$.
\end{proposition}

\begin{proof}[Proof sketch]
One elementary step reads the two neighbour bits of site $i$ and then applies
\rt{I} or $X$ to site $i$. The controls are bits $\neq i$ and are untouched by
that step, so the step is a controlled-$X$ and hence a bijection; a composition of
bijections is a bijection, whatever the layer order. Conversely \rt{D} and \rt{E}
send both values of the target bit to the same value, so two states differing only
in site $i$ but sharing its neighbour pattern collide.
\end{proof}

\noindent Note the proof does not use the sublattice structure, so it survives the
case where the even sites fail to be an independent set (\rt{pbc} at odd $N$).
Checked exhaustively anyway at $N=7,8,9$ for both conventions: exactly those $16$
are bijections and none of the other $240$ is.

\paragraph{Not the six reversible ECA.} A better-known count lives nearby and the
two must not be conflated. For \emph{elementary} cellular automata --- one layer,
simultaneous update, each cell a function of its three-cell neighbourhood ---
only \textbf{six} of the $256$ rules are reversible: $15$, $51$, $85$, $170$,
$204$, $240$, the shifts and complements.

Two things differ, not one. The construction differs: the brick-wall QCA updates
one sublattice at a time and applies a gate to the target conditioned on its
neighbours, which makes invertibility a purely \emph{local} property of the gate.
And the quantifier differs: the ECA six are those bijective for \emph{every} $N$,
an intersection over lattice sizes, and at a fixed $N$ the bijective set is
strictly larger --- on a ring of $N=7$ it also contains $45$, $75$, $89$ and
$150$. Our $16$ need no such intersection: the proposition above is proved
locally and holds at every $N$ separately, which is why the two counts are not
comparable term by term. Only $51$ and $204$ lie in both sets.

\noindent So the $16$ rules that were ``unitary'' in the Hadamard family are here
\emph{reversible} in the sense just defined: every state is recurrent, each
component is a single cycle, and the sector decomposition \emph{is} the orbit
decomposition of a permutation of $\{0,1\}^N$.

\paragraph{Implementation.} The site order and the neighbour conventions are not
reimplemented. The compiled step list comes straight from
\fn{core/cycle.py::\_compile}, so the brick-wall order (even sites ascending,
then odd, each reading the current state) and the boundary handling are
byte-for-byte those of the validated engine; only the local action is swapped.
Nothing in the Hadamard engine is touched, and the two live in separate stores.
This matters because the brick-wall order is \emph{not} two simultaneous
sublattice updates once the even sites stop being an independent set.

\paragraph{Cost.} The whole analysis is one $O(2^N)$ pass over a functional
graph --- a three-colour walk that finds the cycles and, on the way back, each
node's distance to one. No amplitudes, no Tarjan stack, no union-find needed
(though we run union-find anyway as a cross-check, \S\ref{sec:valid}). All $256$ rules at $N=24$
--- sixteen million states each --- take $43$~min in total, against hours per
\emph{rule} for the Hadamard sweep at $N=16$. The map itself is built with numpy,
lifting the loop over $x$ out of Python while keeping the site order from
\fn{\_compile}, which is what makes $N=24$ affordable; the vectorised and scalar
builders are checked against each other in the tests.

\section{Coverage}
\label{sec:sweep}

\begin{center}
\iftable{tab_r10_coverage_obc0.tex}
\end{center}

\noindent Uniform over all $256$ rules. \rt{pbc} is deliberately not run: the ring
has its own caveats and is held for a separate report.

\section{The headline: large sectors, short attractors}
\label{sec:headline}

\begin{center}
\iftable{tab_r10_families_obc0.tex}
\end{center}

\noindent Read the two middle columns against each other. \textbf{$239$ of $256$
rules have exponentially large sectors, and only $24$ have exponentially long
cycles.} For the $80$ V-free rules the count is \emph{zero}: not one of them has a
cycle whose length grows exponentially, while $78$ of them have exponentially
large sectors. The median cyclic fraction --- the share of the basis that is
periodic at all --- is $1.8\times10^{-7}$ for the V+reset family and
$1.5\times10^{-5}$ for the V-free family at $N=24$, with a median transient depth
of $13$ and $12$ steps.

So the permutation circuits are the textbook classical picture: a state falls
through a long transient into a short attractor, and essentially all of the
Hilbert space is transient. The sector structure is large and the recurrent
structure is tiny, and the gap between them is the whole content of the family.

\begin{figure}[htbp]
\centering
\ifgraphic{../../figures/fig_xgate_maps_obc0.pdf}{\textwidth}
\caption{\textbf{X1, the two maps.} Left: base of the sector count against base of
the largest \emph{sector}. Right: the same $x$-axis --- identical by
\eqref{eq:F1}, not merely similar --- against the base of the longest
\emph{cycle}. The dotted curve is $ab=2$; it binds the left panel, because
sectors partition, and is only a reference on the right, because cycles do not.
The collapse of the right panel onto $b=1$ is the headline: marker area grows with
the number of rules in a cell and the two largest cells there hold $110$ and $40$
rules.}
\label{fig:x1}
\end{figure}

\begin{figure}[htbp]
\centering
\ifgraphic{../../figures/fig_xgate_drop_obc0.pdf}{\textwidth}
\caption{\textbf{X2, the drop, per rule.} Left: longest cycle against largest
sector. Every rule lies on or below the diagonal --- a cycle cannot be longer than
the sector containing it --- and the reversible rules are exactly the ones
\emph{on} it, since for them the sector \emph{is} the cycle. Right: the
distribution of the cyclic fraction, log-scaled, by family.}
\label{fig:x2}
\end{figure}

\section{The reversible rules}
\label{sec:reversible}

For the $16$ rules built from \rt{I} and \rt{V} the two decompositions are not
just equal in number: the sectors \emph{are} the cycles, elementwise. So a single
table describes both maps.

\begin{center}
\iftable{tab_r10_reversible_obc0.tex}
\end{center}

\noindent Three anchors are exact and worth stating, because they are the
analogues of R9's:

\begin{itemize}
\item \textbf{Rule $204$} (\rt{IIII}) is the identity: $2^N$ fixed points,
  $a=2$, $b=1$, $ab=2$ --- on the curve, as in the Hadamard case, but for a
  different reason (there every sector was a frozen singleton; here every
  \emph{cycle} is).
\item \textbf{Rule $51$} (\rt{VVVV}) is the global flip $x\mapsto\bar x$:
  $2^{N-1}$ two-cycles, so $a=2$, $b=1$ and $ab=2$. Note the contrast with the
  Hadamard family, where $51$ gives \emph{one} sector of size $2^N$ --- the same
  symbol table, opposite extremes of the sector axis, purely from the gate.
\item \textbf{Rule $150$} (\rt{IVVI}) with $V=X$ is the linear map
  $x_i\mapsto x_i\oplus x_{i-1}\oplus x_{i+1}$ applied brick-wall. Its longest
  cycle is short and grows \emph{linearly}: $D_{\max}=14,15,16,17,18$ at
  $N=13\ldots17$, i.e.\ $N+1$ over the range computed, with
  $n_\wcc\approx2^N/(N+1)$ so $a=2$, $b=1$, $ab=2$. Against the Hadamard rule
  $150$, whose sectors are the binomial domain-wall shells of R8, this is a
  completely different object built from the same table.
\end{itemize}

\noindent The family spans the whole plane. At one end $204$ and $51$ have $2^{N}$
and $2^{N-1}$ tiny orbits; at the other, $57$ and $99$ (\rt{VIVV}, \rt{VVIV})
have only $110$ sectors at $N=17$ with the largest holding $36\,456$ states,
giving $a=1.405$, $b=1.961$ and $ab=2.755$. Four rules --- $60$, $102$, $153$,
$195$ --- and also $108$ and $201$ sit at $b=2$ exactly with $a\approx1.79$--$1.85$,
so $ab\approx3.6$--$3.7$, far above the exclusion curve: many orbits \emph{and} a
largest orbit that is a fixed fraction of the space.

\section{What carries over, and what does not}
\label{sec:carry}

\paragraph{Carries over.} The sum rule and the exclusion curve are properties of a
\emph{partition}, and sectors partition here exactly as they did there. The
fit-free finite-$N$ form
\[
n_\wcc(N)\cdot D_{\max}(N)\;\ge\;2^N
\]
holds for all $256$ rules at every $N$, and is tight in the same place --- the
smallest ratio anywhere is exactly $1.0000$, at rule $0$, $N=6$.

\paragraph{Does not carry over.} Everything that depended on the Hadamard being a
\emph{superposing} gate. There is no interference, so no destructive cancellation
in $\operatorname{succ}$; no dark states; no distinction between the monitored and
the unmonitored channel at the level of the support, since a single Kraus branch
survives on each basis state. The Tier-1d monitored/unmonitored certificate is
vacuous here. Conversely the transient structure, which was a thin and somewhat
awkward part of the Hadamard story, becomes the dominant feature.

\paragraph{The one identity that is a genuine control.} $n_\rec=n_\wcc$ held for
\textbf{every unit computed}, with zero failures, and the union-find sector
decomposition --- which knows nothing about functional graphs --- reproduced the
basin decomposition exactly wherever both were run. Two independent code paths,
one theorem, no discrepancies.

\begin{figure}[htbp]
\centering
\ifgraphic{../../figures/fig_xgate_transient_obc0.pdf}{\textwidth}
\caption{\textbf{X3, transient structure.} Cyclic fraction (log scale), transient
depth, and largest sector as a fraction of the space, against $N$. The reversible
rules sit at cyclic fraction $1$ and depth $0$ by construction; the irreversible
ones fall away exponentially.}
\label{fig:x3}
\end{figure}

\section{Validation}
\label{sec:valid}

\begin{center}
\iftable{tab_r10_validation_obc0.tex}
\end{center}

\noindent The last row is the point of the family rather than a failure: the cycle
map has no exclusion curve, because cycles do not partition the basis, and $217$
of $241$ rules do sit below $ab=2$ there. In the sector map, where the curve does
bind, only \emph{two} rules fall below --- and neither is a violation. Both are
the degenerate case R9~\S5 documents, where the sector count is sub-exponential
so the base product silently drops the polynomial factor: $X_{25}$
(\rt{VIVD}) has an essentially constant sector count
($2,3,2,2,3,2$ over $N=19\ldots24$) with $b=1.99887$, and $X_{182}$ (\rt{IVVE})
a linear one ($11,11,12,12,13,13$) with $b=1.881$. With $a=1$ the inequality
degenerates to $b\ge2$, approached from below at finite $N$, which is what both
are doing. Carrying the sweep from $N=20$ to $N=24$ took this count from $14$ to
$2$, so the residue is a finite-size artefact and not a family of offenders.

\section{Open items}

\begin{enumerate}
\item \textbf{\rt{pbc}.} Held, as for R9.
\item \textbf{Rule $150$'s cycle length.} $D_{\max}=N+1$ over $N=13\ldots17$ is
  five points; the map is $\mathrm{GF}(2)$-linear, so the orbit structure is
  governed by the order of a matrix over $\mathrm{GF}(2)$ and should be derivable
  rather than fitted. That derivation is not attempted here.
\item \textbf{The $24$ rules with exponentially long cycles.} Six are
  reversible (where a long cycle is unsurprising); the other eighteen are
  V+reset rules that nonetheless sustain an exponentially long attractor despite
  a median cyclic fraction of $10^{-7}$. That list is short and deserves the same
  treatment R9 gave its eight open-system-fragmented rules.
\item \textbf{Comparison with the Hadamard sector map} rule by rule. The two
  families share a symbol table but not a single quantitative law; whether any
  rule-level correlation survives the gate change is untested.
\end{enumerate}

\section{Reproduce}

\begin{verbatim}
python -m qca_fragmentation.permutation.sweep --bc obc0 --n-min 6 --n-max 20
python -m qca_fragmentation.permutation.analysis --bc obc0
python -m qca_fragmentation.permutation.tables   --bc obc0
python -m qca_fragmentation.permutation.figures  --bc obc0 --rebuild
python -m pytest qca_fragmentation/tests/test_xgate.py -q
\end{verbatim}

\begin{sloppypar}
Code: \fn{permutation/xca.py} (the map, ported from \fn{core/cycle.py}),
\fn{permutation/functional.py} (the one-pass decomposition and the union-find
cross-check), \fn{permutation/analysis.py}, \fn{permutation/tables.py},
\fn{permutation/figures.py}. Data:
\fn{results\_xgate/\{rule\}\_\{bc\}.jsonl} and
\fn{analytics/xgate\_map\_obc0.json}.
\end{sloppypar}

\end{document}
